% TOP_PROBLEM: {"id": 2737, "title": "Kirby Problem 1.78", "category_id": 7, "category": {"id": 7, "name": "topology", "display_name": "Topology", "description": "Properties preserved under continuous deformations.", "slug": "topology", "order_index": 7, "created_at": "2026-07-31T15:26:25.671Z"}, "status": "open", "problem_number": "KP-1.78", "set_id": 11}
% FIRST_POSTED: 2026-10-01
% ENTRY_KIND: statement_only
% UnsolvedMath ID 2737; source catalogue code KP-1.78
% !TeX program = lualatex
% Definition and sources only; this file is not an LLM solution attempt.
\documentclass[11pt,a4paper]{article}
\usepackage[margin=25mm]{geometry}
\usepackage{fontspec}
\setmainfont{lmroman10-regular.otf}[BoldFont=lmroman10-bold.otf,
 ItalicFont=lmroman10-italic.otf,BoldItalicFont=lmroman10-bolditalic.otf]
\usepackage{amsmath,amssymb,mathtools}
\usepackage{unicode-math}
\setmathfont{latinmodern-math.otf}
\AtBeginDocument{\let\setminus\smallsetminus}
\usepackage{xurl,hyperref}
\hypersetup{colorlinks=true,urlcolor=blue}
\setcounter{secnumdepth}{0}
\setlength{\parindent}{0pt}
\setlength{\parskip}{5pt}
\setlength{\emergencystretch}{3em}
\begin{document}
\section*{Kirby Problem 1.78}\label{2737}
{\small UnsolvedMath ID 2737 · KP-1.78 · Topology · Kirby's Problems in Low-Dimensional Topology}
\subsection{Definitions and mathematical statement}
Let $D$ be a diagram of the unknot and let $c(D)$ count its crossings.
Among all finite sequences of Reidemeister moves from $D$ to a zero-crossing
circle, define the smallest achievable maximum crossing number by
\[
h(D)=\min_{D=D_0\to D_1\to\cdots\to D_m=U}
          \max_{0\leq i\leq m}c(D_i).
\]
Planar isotopy is allowed freely between moves. All intermediate diagrams
represent the unknot, but their crossing numbers can increase.
For each integer $n\geq0$, put
\[
h(n)=\max\{h(D):D\text{ is an unknot diagram with exactly }n\text{ crossings}\}.
\]
Use combinatorial plane diagram types modulo planar isotopy, making the outer
maximum finite. The inner minimum exists because the Reidemeister moves connect
all diagrams of the unknot.

Determine $h(n)$, or its sharp growth as $n\to\infty$. Since the initial diagram
is included, $h(n)\geq n$. A hard unknot diagram satisfies $h(D)>c(D)$, meaning
every simplification must first pass above its original crossing count. The
problem concerns this necessary crossing excess, independent of how many moves
the chosen simplification uses.
\subsection{Short English statement}
How many extra crossings can be unavoidable while simplifying an $n$-crossing diagram of the unknot?
\subsection{Sources}
\begin{itemize}
\item[S1] ULAM.AI and the UnsolvedMath Contributors, supplied problems.json, record 2737 (KP-1.78), Kirby's Problems in Low-Dimensional Topology. Catalogue identity and source transcription; CC BY 4.0.
\url{https://huggingface.co/datasets/ulamai/UnsolvedMath}.
\item[S2] K3: A New Problem List in Low-Dimensional Topology, Problem 1.78. Expanded formulation of the supplied catalogue statement.
\url{https://math.berkeley.edu/sites/default/files/surv-295-ruberman-watermarked-author-pdf.pdf}.
\end{itemize}
\end{document}
